\documentclass[12pt,a4paper,twoside,openright]{book}

%% Pacotes
\usepackage[utf8]{inputenc}
\usepackage[T1]{fontenc}
\usepackage[english]{babel}
\usepackage{lmodern}
\usepackage{geometry}
\usepackage{amsmath,amsthm,amssymb}
\usepackage{hyperref}
\usepackage{xcolor}
\usepackage{enumitem}
\usepackage{microtype}
\usepackage{fancyhdr}
\usepackage{titlesec}
\usepackage{tocloft}
\usepackage{booktabs}
\usepackage{array}
\usepackage{mdframed}
\usepackage{cleveref}

\geometry{
  a4paper,
  left=3.2cm, right=2.8cm,
  top=3cm, bottom=3cm,
  headheight=15pt
}

%% Cores
\definecolor{darkblue}{HTML}{1a1a2e}
\definecolor{accent}{HTML}{c0392b}
\definecolor{linkcolor}{HTML}{2980b9}
\definecolor{thmback}{HTML}{f0f4ff}
\definecolor{conjback}{HTML}{fff5f5}
\definecolor{defback}{HTML}{f5fff5}

%% Hyperref
\hypersetup{
  colorlinks=true,
  linkcolor=darkblue,
  urlcolor=linkcolor,
  citecolor=accent,
  pdftitle={Cobertura e Partição em Estruturas Monocromáticas},
  pdfauthor={Survey compilado para qualificação},
  bookmarksnumbered=true
}

%% Ambientes de teoremas
\theoremstyle{definition}
\newmdtheoremenv[
  backgroundcolor=thmback,
  linecolor=darkblue,
  linewidth=1.5pt,
  leftmargin=0pt, rightmargin=0pt,
  innertopmargin=6pt, innerbottommargin=6pt,
  innerleftmargin=10pt, innerrightmargin=10pt,
  skipabove=8pt, skipbelow=6pt
]{theorem}{Theorem}[chapter]

\newmdtheoremenv[
  backgroundcolor=conjback,
  linecolor=accent,
  linewidth=1.5pt,
  leftmargin=0pt, rightmargin=0pt,
  innertopmargin=6pt, innerbottommargin=6pt,
  innerleftmargin=10pt, innerrightmargin=10pt,
  skipabove=8pt, skipbelow=6pt
]{conjecture}[theorem]{Conjecture}

\newmdtheoremenv[
  backgroundcolor=defback,
  linecolor=darkblue!60,
  linewidth=1pt,
  leftmargin=0pt, rightmargin=0pt,
  innertopmargin=5pt, innerbottommargin=5pt,
  innerleftmargin=10pt, innerrightmargin=10pt,
  skipabove=6pt, skipbelow=4pt
]{definition}[theorem]{Definition}

\newtheorem{corollary}[theorem]{Corollary}
\newtheorem{lemma}[theorem]{Lemma}
\newtheorem{proposition}[theorem]{Proposition}
\newtheorem{remark}[theorem]{Remark}
\newtheorem{problem}{Open Problem}[chapter]

%% Título das seções em darkblue
\titleformat{\chapter}[display]
  {\normalfont\huge\bfseries\color{darkblue}}
  {\chaptertitlename\ \thechapter}{20pt}{\Huge}

\titleformat{\section}
  {\normalfont\Large\bfseries\color{darkblue}}
  {\thesection}{1em}{}

\titleformat{\subsection}
  {\normalfont\large\bfseries\color{darkblue!80}}
  {\thesubsection}{1em}{}

%% Cabeçalho
\pagestyle{fancy}
\fancyhf{}
\fancyhead[LE]{\small\color{darkblue}\leftmark}
\fancyhead[RO]{\small\color{darkblue}\rightmark}
\fancyfoot[C]{\thepage}

%% Macros úteis
\newcommand{\arxiv}[1]{\href{https://arxiv.org/abs/#1}{\textcolor{linkcolor}{arXiv:#1}}}
\newcommand{\Kn}{K_n}
\newcommand{\Knm}{K_{n,n}}
\newcommand{\Knnn}{K_{n,n,n}}
\newcommand{\Gnp}{G(n,p)}
\newcommand{\dmin}{\delta(G)}
\newcommand{\cpr}{\mathrm{cp}_r}

\begin{document}

%%===================================================================
%% CAPA
%%===================================================================
\begin{titlepage}
\centering
\vspace*{2cm}
{\color{darkblue}\rule{\linewidth}{2pt}}\\[0.5cm]
{\Huge\bfseries\color{darkblue}
  Cobertura e Partição\\[4pt]
  em Estruturas\\[4pt]
  Monocromáticas}\\[0.8cm]
{\color{accent}\rule{\linewidth}{1pt}}\\[0.6cm]
{\large
  Caminhos $\cdot$ Ciclos $\cdot$ Árvores $\cdot$ Componentes Conexas}\\[0.3cm]
{\normalsize
  Grafos Completos $\cdot$ Grau Mínimo Alto $\cdot$ Aleatórios $\cdot$ Bipartidos $\cdot$ Multipartidos}\\[2cm]
{\large\textbf{Versão I} — Enunciados e Contexto}\\[0.4cm]
{\normalsize Survey compilado a partir de 44 artigos do arXiv}\\[0.3cm]
{\normalsize Para uso em qualificação de doutorado}\\[2.5cm]
{\normalsize Compilado em julho de 2026}
\vfill
{\color{darkblue}\rule{\linewidth}{2pt}}
\end{titlepage}

\frontmatter
\tableofcontents
\cleardoublepage

%%===================================================================
\chapter*{Prefácio / Preface}
\addcontentsline{toc}{chapter}{Prefácio / Preface}
%%===================================================================


Este documento é um survey de 44 artigos do arXiv sobre um tema central da combinatória extremal:
\textbf{a existência de estruturas monocromáticas que cobrem ou particionam os vértices de grafos com coloração de arestas}.

As estruturas de interesse são \textbf{caminhos}, \textbf{ciclos}, \textbf{árvores} e \textbf{componentes conexas} monocromáticas.
Os grafos considerados variam entre grafos completos, grafos com grau mínimo alto, grafos aleatórios,
grafos bipartidos e grafos multipartidos.

\medskip
\noindent\textbf{Como usar este documento.}
Cada capítulo corresponde a uma categoria de resultados. Dentro de cada capítulo, os artigos são
apresentados em ordem cronológica ou temática, com os enunciados exatos dos principais teoremas
(em inglês, conforme aparecem nos artigos originais), seguidos de contexto e discussão em português.

\medskip
\noindent\textbf{Convenção.} Ao longo do texto, um \emph{ciclo} pode ser o conjunto vazio, um vértice isolado
ou uma aresta (conforme convenção padrão nesta área). O mesmo vale para caminhos.

\medskip
\noindent\textbf{Referências.} Cada seção inclui o identificador arXiv do artigo correspondente.
Todos os PDFs estão disponíveis na raiz do repositório.


\mainmatter

%%===================================================================
\chapter{Introdução e Definições}
%%===================================================================


\section{Contexto histórico}

O estudo de estruturas monocromáticas em grafos com coloração de arestas remonta ao clássico Teorema de Ramsey
(1930), que garante a existência de subgrafos monocromáticos em qualquer coloração de grafos completos suficientemente grandes.
Uma variação natural e rica desse problema é a seguinte: em vez de encontrar \emph{uma} estrutura monocromática,
queremos encontrar poucas estruturas que juntas \emph{cubram ou particionem} todos os vértices.

A pedra fundadora desta área é um resultado de Gerencsér e Gyárfás (1967):


\begin{theorem}[Gerencsér--Gyárfás, 1967]
In every $2$-edge-colouring of $\Kn$, the vertices can be covered by two vertex-disjoint monochromatic paths.
\end{theorem}


\noindent Esse resultado motivou uma série de generalizações e conjecturas que estruturam todo o campo.
As três mais influentes são:


\begin{conjecture}[Gyárfás, 1977]\label{conj:gyarfas-paths}
The vertices of every $r$-edge-coloured complete graph can be covered by $r$ vertex-disjoint monochromatic paths.
\end{conjecture}

\begin{conjecture}[Erdős--Gyárfás--Pyber, 1991]\label{conj:egp-cycles}
The vertices of every $r$-edge-coloured complete graph can be partitioned into $r$ vertex-disjoint monochromatic cycles.
\end{conjecture}

\begin{conjecture}[Ryser, 1967 / Lovász, 1975 — versão para grafos]\label{conj:ryser}
For every $r$-edge-coloured graph $G$, the vertices of $G$ can be covered by at most $(r-1)\alpha(G)$
monochromatic trees, where $\alpha(G)$ is the independence number.
\end{conjecture}


\section{Definições essenciais}


\begin{definition}
Let $G$ be an $r$-edge-coloured graph. We define:
\begin{itemize}[leftmargin=*]
  \item $\mathrm{cp}_r(G)$: the minimum $k$ such that every $r$-colouring of $G$ admits a vertex-partition into
        $k$ monochromatic cycles.
  \item $\mathrm{pp}_r(G)$: the minimum $k$ for vertex-partition into $k$ monochromatic paths.
  \item $\mathrm{tp}_r(G)$: the minimum $k$ for vertex-partition into $k$ monochromatic trees (connected subgraphs).
  \item $\mathrm{tc}_r(G)$: the minimum $k$ for vertex-\emph{cover} (not necessarily partition) by $k$ monochromatic
        connected subgraphs.
\end{itemize}
\end{definition}


\section{Técnicas recorrentes}

Os principais métodos utilizados nos papers desta compilação são:
\begin{itemize}[leftmargin=*]
  \item \textbf{Regularity Lemma de Szemerédi}: reduz o problema a grafos regularizados; técnica central
        em resultados assintóticos dos anos 1990--2010.
  \item \textbf{Método de matchings conexos (Łuczak)}: reduz partição em caminhos/ciclos à existência de
        matchings conexos em grafos de clusters.
  \item \textbf{Método de absorção}: constrói um subgrafo ``absorvente'' que incorpora os vértices restantes;
        essencial em provas exatas.
  \item \textbf{Blow-up Lemma}: permite encontrar subgrafos em grafos regulares; usado em resultados de grau mínimo.
  \item \textbf{Probabilistic method / Lovász Local Lemma}: central nos resultados para grafos aleatórios.
\end{itemize}


%%===================================================================
\chapter{Grafos Completos: Partição em Caminhos e Ciclos}
%%===================================================================


Este capítulo cobre resultados onde o grafo hospedeiro é um grafo completo $K_n$.


\section{A Conjectura de Lehel e sua resolução}

\begin{theorem}[Łuczak--Rödl--Szemerédi, 1998; Allen, 2008; Bessy--Thomassé, 2010]
The vertex set of every $2$-edge-coloured complete graph $\Kn$ can be partitioned into
two monochromatic cycles of \emph{distinct colours}.
\end{theorem}


\noindent\textbf{Nota histórica.} Esta afirmação, conhecida como \emph{Conjectura de Lehel}, foi proposta nos anos 1970.
Foi provada para $n$ grande por Łuczak, Rödl e Szemerédi usando o Regularity Lemma,
depois por Allen sem o Regularity Lemma, e finalmente para todo $n$ por Bessy e Thomassé com uma prova notavelmente curta e elegante.


\section{A falsidade da Conjectura de Erdős--Gyárfás--Pyber para $r\ge3$}

\begin{theorem}[Pokrovskiy, 2014 — \arxiv{1205.5492}]
For $r \geq 3$, there exist infinitely many $r$-edge-coloured complete graphs which cannot be
vertex-partitioned into $r$ monochromatic cycles.
\end{theorem}


\noindent Ou seja, a Conjectura~\ref{conj:egp-cycles} é \textbf{falsa} para $r\ge3$.
O mesmo artigo prova um resultado positivo para 3 cores em termos de caminhos:


\begin{theorem}[Pokrovskiy, 2014 — \arxiv{1205.5492}]
In any $3$-edge-colouring of $\Kn$, the vertices can be covered by three vertex-disjoint monochromatic paths.
\end{theorem}

\begin{remark}
This settles \Cref{conj:gyarfas-paths} for $r=3$. As an intermediate step, the paper shows that in any
$2$-edge-colouring of $\Kn$, the vertices can be covered by a red path and a disjoint blue balanced complete bipartite graph.
\end{remark}

\section{Coberturas com $O(r\log r)$ ciclos}

\begin{theorem}[Gyárfás--Ruszinkó--Sárközy--Szemerédi, 2006]
For large $n$, the vertices of any $r$-edge-coloured $\Kn$ can be partitioned into at most $100r\log r$
vertex-disjoint monochromatic cycles.
\end{theorem}

\section{Coberturas com caminhos da mesma cor: a Conjectura de Erdős--Gyárfás}

\begin{theorem}[Erdős--Gyárfás, 1995]
In every $2$-edge-coloured $\Kn$, there exists a collection of $2\sqrt{n}$ monochromatic paths,
all of the same colour, covering all vertices.
\end{theorem}

\begin{conjecture}[Erdős--Gyárfás, 1995 — resolvida em 2024]
In every $2$-edge-coloured $\Kn$, there exists a collection of $\sqrt{n}$ monochromatic paths,
all of the same colour, covering all vertices.
\end{conjecture}

\begin{theorem}[Pokrovskiy--Versteegen--Williams, 2024 — \arxiv{2409.03623}]
For all $n > 2^{40}$, the vertex set of every $2$-edge-coloured $\Kn$ can be covered by
$\sqrt{n}$ monochromatic paths, all of the same colour.
\end{theorem}


\noindent\textbf{Discussão.} O bound $\sqrt{n}$ é tight: tome uma partição de $V(\Kn)$ em
$A$ e $B$ com $|A|=n-\sqrt{n}+1$, $|B|=\sqrt{n}-1$, colora arestas em $A$ de azul e as demais de vermelho.
Cada caminho vermelho cobre no máximo $2$ vértices de $B$, logo são necessários $\ge \sqrt{n}-1$ caminhos vermelhos para cobrir $B$.


\section{Componentes monocromáticas com muitas arestas}

\begin{theorem}[Luo, 2021 — \arxiv{2111.13342}]
Every $3$-edge-colouring of $\Kn$ contains a monochromatic connected subgraph spanning at least
$\frac{1}{6}$ of the edges.
\end{theorem}

\begin{theorem}[Conlon--Luo--Tyomkyn, 2022 — \arxiv{2204.11360}]
Every $4$-edge-colouring of $\Kn$ contains a monochromatic connected subgraph spanning at least
$\frac{1}{12}$ of the edges.
\end{theorem}

\section{Partição em subgrafos monocromáticos de grau limitado}


Grinshpun e Sárközy (\arxiv{1405.7507}) estudam a partição de grafos completamente coloridos em subgrafos monocromáticos com grau máximo limitado, generalizando os resultados para caminhos e árvores.


%%===================================================================
\chapter{Grafos Completos: Componentes e Árvores}
%%===================================================================

\section{O Teorema de Gyárfás sobre componentes grandes}

\begin{theorem}[Gyárfás, 1977]
In every $r$-edge-colouring of $\Kn$, there is a monochromatic connected component with at least
$\frac{n}{r-1}$ vertices.
\end{theorem}

\section{A Conjectura de Erdős--Gyárfás--Pyber para partição em árvores}

\begin{conjecture}[Erdős--Gyárfás--Pyber, 1991 — parcialmente resolvida]
$\mathrm{tp}_r(K_n) = r-1$, i.e., the vertices of every $r$-edge-coloured $K_n$ can be partitioned into
$r-1$ monochromatic trees.
\end{conjecture}

\begin{theorem}[Haxell--Kohayakawa, 1995]
For $n \geq (1-1/r)^{3(r-1)}$, $\mathrm{tp}_r(K_n) \leq r$.
\end{theorem}


\noindent Ou seja, $r$ árvores são suficientes (em vez das $r-1$ conjecturadas) para $n$ suficientemente grande.
O resultado exato $r-1$ é conhecido apenas para $r \le 5$.


\section{Generalização para hipergrafos completos}

\begin{theorem}[Lichev--Luo, 2023 — \arxiv{2302.04487}]
For a wide range of notions of connectivity in $r$-uniform hypergraphs, the size of the largest
monochromatic component in any $r$-edge-colouring of the complete $r$-uniform hypergraph
$K_n^{(r)}$ is $\Theta(n/(r-1))$, tight up to a factor of $1+o(1)$ as $r$ grows.
\end{theorem}

\begin{theorem}[Bal--DeBiasio, 2023 — \arxiv{2302.06669}]
The Gyárfás theorem on large monochromatic components extends to the class of \emph{expansive hypergraphs}.
\end{theorem}

%%===================================================================
\chapter{Grau Mínimo Alto}
%%===================================================================


Neste capítulo, o grafo hospedeiro $G$ não é necessariamente completo, mas tem grau mínimo $\dmin$ próximo de $n$.


\section{A Conjectura de Balogh--Barát--Gerbner--Gyárfás--Sárközy e sua resolução}

\begin{conjecture}[Balogh--Barát--Gerbner--Gyárfás--Sárközy, 2014 — \arxiv{1509.05544}]
Let $G$ be a graph on $n$ vertices with $\delta(G) \geq \frac{3n}{4}$.
Then for every $2$-edge-colouring of $G$, the vertex set $V(G)$ can be partitioned into two
vertex-disjoint monochromatic cycles, one of each colour.
\end{conjecture}

\begin{theorem}[Letzter, 2015 — \arxiv{1502.07736}]
There exists $n_0$ such that if $G$ is a graph on $n \geq n_0$ vertices with $\delta(G) \geq \frac{3n}{4}$
and its edges are $2$-coloured, then $V(G)$ can be partitioned into two monochromatic cycles of different colours.
\end{theorem}


\noindent\textbf{Técnicas.} A prova usa o Regularity Lemma, o método de matchings conexos de Łuczak,
e o método de absorção. O threshold $3n/4$ é sharp: existem exemplos com $\dmin = \lceil 3n/4 \rceil - 1$
sem tal partição.


\begin{theorem}[2022 — \arxiv{2204.00496}]
A $2$-edge-coloured graph on $n$ vertices with $\delta(G) \geq \frac{2n}{3} + o(n)$ can be partitioned
into \emph{three} monochromatic cycles.
\end{theorem}

\section{Condição de grau mínimo para $r$ cores}

\begin{theorem}[Allen--Böttcher--Lang--Skokan--Stein, 2019 — \arxiv{1902.05882}]
For $r \geq 2$ and $n$ sufficiently large, any $r$-edge-coloured graph $G$ on $n$ vertices with
$\delta(G) \geq \frac{n}{2} + 1200r\log n$ admits a partition into $10^7 r^2$ monochromatic cycles.
\end{theorem}


\noindent A construção que mostra que o número de ciclos é essencialmente ótimo também é dada no artigo.


\section{Condições tipo Ore e Pósa}

\begin{theorem}[Arras, 2022 — \arxiv{2202.06388}]
If $\deg(u)+\deg(v) \geq \frac{4n}{3}+o(n)$ for all non-adjacent $u,v \in V(G)$, then all but
$o(n)$ vertices can be partitioned into two monochromatic cycles of distinct colours (approximate Ore-type version).
\end{theorem}

\section{Resultado geral para $r$ cores e grau $(1-\delta)n$}

\begin{theorem}[Di Braccio--Patel, 2026 — \arxiv{2601.22117}]\label{thm:dibraccio}
There exist constants $k, K > 0$ such that for all $r \geq 2$ and $\delta \in (0, 1/2)$,
\[
  kr\left\lceil\frac{r}{\log(1/\delta)}\right\rceil
  \;\leq\; \cpr(\delta)
  \;\leq\; Kr\log r \left\lceil\frac{r}{\log(1/\delta)}\right\rceil,
\]
where $\cpr(\delta)$ is the minimum number of cycles needed to cover any $r$-coloured $n$-vertex graph
with $\delta(G) \geq (1-\delta)n$.
\end{theorem}


\noindent\textbf{Significado.} Este resultado de 2026 resolve (a menos de um fator $\log r$) a questão levantada
independentemente por Letzter e Pokrovskiy. Também \textbf{disprova uma conjectura de Bal e DeBiasio}
sobre coberturas com árvores monocromáticas.


\section{Componentes monocromáticas e grau mínimo}

\begin{theorem}[Guggiari--Scott, 2019 — \arxiv{1909.09178}]
The maximum $\varepsilon$ such that every $3$-edge-colouring of an $n$-vertex graph with
$\delta(G) \geq (1-\varepsilon)n$ contains a monochromatic component of order $\geq \frac{n}{2}$
is $\varepsilon_3 = \frac{1}{6}$.
\end{theorem}


\noindent Este resultado \textbf{disprova uma conjectura de Gyárfás e Sárközy} que sugeria um valor maior para $\varepsilon_3$.


\begin{theorem}[Bal--DeBiasio, 2023 — \arxiv{2301.05806}]
In every $(k+1)$-edge-colouring of a $k$-uniform hypergraph on $n$ vertices with minimum
$(k-1)$-degree $\geq \frac{kn}{k+1} - (k-1)$, there is a monochromatic component of order
$\geq \frac{kn}{k+1}$. This bound on the codegree is best possible.
\end{theorem}

%%===================================================================
\chapter{Grafos Aleatórios}
%%===================================================================


Neste capítulo, o grafo hospedeiro é o grafo aleatório binomial $G(n,p)$.
O objetivo é determinar, em função de $p$ e $r$, quantos ciclos/árvores/componentes monocromáticos
são necessários para cobrir os vértices com alta probabilidade (c.a.s.\ / w.h.p.).


\section{Partição em componentes: extensão de Haxell--Kohayakawa}

\begin{theorem}[Bal--DeBiasio, 2015 — \arxiv{1509.09168}]
If $p \geq \left(\frac{27\log n}{n}\right)^{1/3}$, then w.h.p.\ every $2$-edge-colouring of $G(n,p)$
admits a vertex-partition into two monochromatic components.
Moreover, if $p = \omega(1)/n$, then w.h.p.\ every $r$-colouring of $G(n,p)$ contains a monochromatic
component of order $\geq \left(1-o(1)\right)\frac{n}{r-1}$.
\end{theorem}

\section{Partição em ciclos: análogo aleatório de Erdős--Gyárfás--Pyber}

\begin{theorem}[Korándi--Sárközy--Selkow, 2018 — \arxiv{1807.06607}]
Let $r \geq 2$ and $p = p(n) \geq 29 r^5 n^{-1/(2r)}$. Then w.h.p.\ any $r$-edge-colouring of
$G \sim G(n,p)$ admits a partition into at most $1000r^4\log r$ monochromatic cycles.
\end{theorem}

\begin{remark}
A construction of Bal and DeBiasio shows that $p = o\!\left((r\log n/n)^{1/r}\right)$ is insufficient.
The exact threshold is conjectured to be of order $(r\log n/n)^{1/r}$.
\end{remark}

\begin{theorem}[Korándi--Sudakov, 2017 — \arxiv{1712.03145}]
For appropriate edge density $p$, any $r$-edge-colouring of $G(n,p)$ can be covered by
$O(r^8\log r)$ vertex-disjoint monochromatic cycles.
\end{theorem}

\section{Árvores monocromáticas em grafos aleatórios}

\begin{theorem}[Kohayakawa--Mota--Schacht, 2016 — \arxiv{1611.10299}]
The threshold for the $2$-colour Ramsey tree-partition property
(i.e., w.h.p.\ every $2$-colouring of $G(n,p)$ has a vertex-partition into two monochromatic trees of distinct colours)
is determined.
\end{theorem}

\begin{theorem}[Bradač--Bucić, 2021 — \arxiv{2109.02569}]
The problem of covering $G(n,p)$ with monochromatic components is connected to
a Helly-type question for hypergraphs raised by Erdős, Hajnal, and Tuza.
A modified version of this hypergraph problem controls the covering number more precisely.
\end{theorem}

\section{Hipergrafos aleatórios}

\begin{theorem}[Bennett--DeBiasio--Dudek--English, 2017 — \arxiv{1709.02990}]
The results of Gyárfás and Füredi on large monochromatic components in colourings of complete
$k$-uniform hypergraphs extend to the random hypergraph setting. Long monochromatic loose cycles
in $r$-coloured random hypergraphs are also studied.
\end{theorem}

\section{Componentes com muitas arestas em grafos aleatórios}

\begin{theorem}[Fox--Luo, 2025 — \arxiv{2509.01766}]
Results on the maximum number of edges in a monochromatic component in $r$-colourings of $K_n$
(Luo for $r=3$; Conlon--Luo--Tyomkyn for $r=4$) extend to sparse random graphs and to graphs
of large minimum degree.
\end{theorem}

%%===================================================================
\chapter{Grafos Bipartidos}
%%===================================================================

\section{Caminhos longos em $K_{n,n}$: a Conjectura de Gyárfás--Lehel}

\begin{theorem}[Gyárfás--Sárközy, 2018 — \arxiv{1806.05119}]
In every $2$-edge-colouring of $K_{n,n}$, there exists a monochromatic path on at least $2\lceil n/2\rceil$ vertices.
This bound is tight.
\end{theorem}

\begin{theorem}[Gyárfás--Sárközy, 2018 — \arxiv{1806.05119}]
If $G$ is a balanced bipartite graph on $2n$ vertices with $\delta(G) \geq (3/4+o(1))n$, then in every
$2$-edge-colouring of $G$ there exists a monochromatic cycle on at least $(1+o(1))n$ vertices,
or the colouring is close to extremal.
\end{theorem}

\section{Partição de $K_{n,n}$ em caminhos}

\begin{theorem}[Pokrovskiy, 2014]
Let $G = K_{n,n}$ with edges $2$-coloured. If the colouring is \emph{not} a split colouring, then
the vertices of $G$ can be partitioned into two monochromatic paths of different colours.
\end{theorem}


\noindent\textbf{Split colouring:} uma 2-coloração onde todas as arestas de um lado têm a mesma cor.


\section{Partição em caminhos e ciclos: generalização para hipergrafos}

\begin{theorem}[Stein, 2022 — \arxiv{2204.12464}]
\textbf{(a)} Any $2$-edge-coloured complete $3$-uniform hypergraph $K_n^{(3)}$ can be partitioned into
two monochromatic vertex-disjoint tight paths of different colours.\\
\textbf{(b)} Any $2$-edge-coloured $K_{n,n}$ that is not a split colouring can be partitioned into
a monochromatic path and a monochromatic cycle of different colours.
\end{theorem}

\begin{conjecture}[Stein, 2022 — \arxiv{2204.12464}]
For any $r, n \in \mathbb{N}$, every $2$-edge-colouring of the $r$-uniform complete hypergraph $K_n^{(r)}$
has a partition into two monochromatic vertex-disjoint tight paths of different colours.
\end{conjecture}

\section{Partição de $K_{n,n}$ em ciclos}

\begin{theorem}[2024 — \arxiv{2409.03394}]
For every $2$-edge-colouring of the complete balanced bipartite graph $K_{n,n}$, the vertices can be
partitioned into at most $4$ monochromatic cycles.
\end{theorem}


\noindent\textbf{Nota.} O bound $4$ é esperado não ser ótimo; a conjectura é que $3$ ciclos bastam para $n$ grande.
O limite inferior é $2r-1$ para $r$ cores gerais (construção simples por extensão dos resultados de $K_{n,n}$).


\section{Coberturas em bipartidos aleatórios e de grau mínimo alto}

\begin{theorem}[2024 — \arxiv{2403.12587}]
\textbf{(a) [Random]} For $p \gg ({\log n}/n)^{1/2}$, w.h.p.\ every $2$-colouring of $G(n,n,p)$
admits a cover of all but $O(1/p)$ vertices using at most $3$ vertex-disjoint monochromatic components.\\
\textbf{(b) [Minimum degree]} Every $2$-colouring of a bipartite graph with parts of size $n$ and
$\delta(G) \geq \left(\frac{13}{16}+o(1)\right)n$ admits a cover by at most $3$ monochromatic components.
\end{theorem}

\section{Colorações locais em bipartidos completos}

\begin{theorem}[Gyárfás--Sárközy, 2015 — \arxiv{1501.05619}]
For any $2$-local colouring of $K_{n,n}$
(i.e., every vertex is incident with at most $2$ colours), the vertices can be covered by at most
$3$ vertex-disjoint monochromatic paths.
\end{theorem}

\section{Componentes grandes em bipartidos}

\begin{theorem}[DeBiasio--Krueger--Sárközy, 2018 — \arxiv{1806.05271}]
In every $r$-edge-colouring of $K_{m,n}$, there is a monochromatic connected component with at least
$\frac{m+n}{r}$ vertices.
\end{theorem}

\begin{conjecture}[DeBiasio--Krueger--Sárközy, 2018 — \arxiv{1806.05271}]\label{conj:bipartite-comp}
In every $r$-edge-colouring of a bipartite graph $G$ with parts $X$ and $Y$ ($|X|=m$, $|Y|=n$),
$\delta(X,Y) > \left(1-\frac{1}{r+1}\right)n$ and $\delta(Y,X) > \left(1-\frac{1}{r+1}\right)m$,
there is a monochromatic component on $\geq \frac{m+n}{r}$ vertices.
(Proved for $r=2$; weaker bound for $r \geq 3$.)
\end{conjecture}

\begin{theorem}[Füredi--Luo, 2020 — \arxiv{2008.12217}]
``Almost complete'' graphs (i.e., $K_n$ minus few edges) and almost complete bipartite graphs retain the Gyárfás property:
every $(t+1)$-edge-colouring contains a monochromatic component of order $\geq n/t$.
\end{theorem}

\section{Ciclos em bipartidos com grau mínimo alto}

\begin{theorem}[Zhang--Peng--Luo, 2023 — \arxiv{2304.08003}]
For any $m > n$, if $G$ is a balanced bipartite graph of order $2(m+n-1)$ with
$\delta(G) > \frac{3}{4}(m+n-1)$, then $G \to (C_m^M, C_n^M)$, where $C_i^M$ denotes a
monochromatic cycle of length $i$.
\end{theorem}

%%===================================================================
\chapter{Grafos Multipartidos}
%%===================================================================

\section{Cobertura de grafos $k$-partidos com 2 caminhos}

\begin{theorem}[Schaudt--Stein, 2014 — \arxiv{1407.5083}]\label{thm:multipartite-paths}
Let $G$ be a complete $k$-partite graph on $n$ vertices, $k \geq 3$, whose edges are $2$-coloured.
If the largest partition class contains at most $n/2$ vertices (\emph{fair} colouring), then the vertices
of $G$ can be covered by two vertex-disjoint monochromatic paths of distinct colours.
\end{theorem}

\begin{theorem}[Schaudt--Stein, 2014 — \arxiv{1407.5083}]
Under the same hypotheses, all but $o(n)$ vertices can be covered by two monochromatic cycles,
provided the colouring is not close to a split colouring.
Consequently, for large enough $n$, the whole graph can be partitioned into at most $14$ monochromatic cycles.
\end{theorem}

\section{Caminhos e ciclos longos em multipartidos}

\begin{theorem}[Balogh--Kostochka--Lavrov--Liu, 2019 — \arxiv{1905.04657}]
For fixed $s$ and large $n$, and for every $2$-edge-colouring of $K_{n_1,\ldots,n_s}$, exact conditions
on $n_1,\ldots,n_s$ are determined for the existence of:
\begin{itemize}
  \item a monochromatic even cycle $C_{2n}$;
  \item a monochromatic cycle $C_{\geq 2n}$;
  \item a monochromatic path $P_{2n}$;
  \item a monochromatic path $P_{2n+1}$.
\end{itemize}
In particular, this confirms for large $n$ the Gyárfás--Ruszinkó--Sárközy--Szemerédi conjecture that
in every $2$-colouring of $\Knnn$ there is a monochromatic path $P_{2n+1}$.
\end{theorem}

\begin{theorem}[Balogh--Kostochka--Lavrov--Liu, 2019 — \arxiv{1905.04653}]
Exact Ramsey-type bounds on the sizes of monochromatic connected matchings in $2$-edge-coloured
complete multipartite graphs, together with a stability theorem.
\end{theorem}


\noindent\textbf{Nota.} Um \emph{matching conexo} é um matching cujas arestas pertencem todas à mesma componente conexa do grafo.
Este conceito, introduzido por Łuczak, é a ferramenta central para reduzir problemas de caminhos/ciclos em grafos regulares.


\section{Hipergrafos $r$-partidos: cobertura com poucos componentes}

\begin{theorem}[Hawranick--Luo, 2026 — \arxiv{2603.04704}]
For $k \geq 2r \geq 6$, in any spanning $k$-edge-colouring of a complete $r$-partite $r$-uniform hypergraph $H$,
the vertices of $H$ can be covered by at most $k-r+1$ monochromatic connected components.
Moreover, for $k \in \{2,3\}$, every spanning $k$-edge-colouring of a complete bipartite graph admits
a covering using at most $k$ monochromatic components.
\end{theorem}


\noindent\textbf{Significado.} Este resultado de 2026 prova uma conjectura de Gyárfás e Király, que é um caso especial
da \emph{Conjectura de Ryser} (\Cref{conj:ryser}).


\section{Hipergrafos partidos: cobertura clássica}

\begin{theorem}[Gyárfás--Király, 2016 — \arxiv{1604.02791}]
In every $k$-colouring of a complete $r$-uniform $r$-partite hypergraph, the vertices can be covered
by at most $\lceil k/r \rceil$ monochromatic connected components.
\end{theorem}

%%===================================================================
\chapter{Problemas Abertos e Conjecturas}
%%===================================================================


Este capítulo compila os principais problemas em aberto identificados nos 44 artigos desta compilação.
São organizados por dificuldade e relevância para a qualificação.


\begin{problem}[Conjectura de Gyárfás — Caminhos em $r$ cores]
Prove \Cref{conj:gyarfas-paths}: the vertices of every $r$-edge-coloured $K_n$ can be covered by
$r$ vertex-disjoint monochromatic paths.

\medskip\noindent
\textit{Status:} Proved for $r=2$ (Gerencsér--Gyárfás), $r=3$ (Pokrovskiy 2014). Open for $r \geq 4$.
\end{problem}

\begin{problem}[Conjectura de Erdős--Gyárfás--Pyber — Ciclos, versão enfraquecida]
Is $\mathrm{cp}_r(K_n) = O(r)$? The current best bound is $O(r\log r)$.
\end{problem}

\begin{problem}[Conjectura de Pokrovskiy — 3 ciclos com $\delta \geq 2n/3$]
Show that every $2$-edge-coloured $n$-vertex graph with $\delta(G) \geq \frac{2n}{3}$ can be exactly
partitioned into $3$ monochromatic cycles.

\medskip\noindent
\textit{Status:} Proved approximately by Allen et al.; the exact version is open.
\end{problem}

\begin{problem}[Conjectura de Bal--DeBiasio — Threshold para árvores em $G(n,p)$]
Is the threshold for $\mathrm{tp}_r(G(n,p)) \leq r$ equal to $p = (1+\varepsilon)(r\log n/n)^{1/r}$?

\medskip\noindent
\textit{Status:} Confirmed for $r=2$ (Kohayakawa--Mota--Schacht). Open for $r \geq 3$.
Note: Di Braccio--Patel 2026 disproved a related tree-covering conjecture.
\end{problem}

\begin{problem}[Conjectura de Gyárfás--Lehel — Componentes em bipartidos]
Prove \Cref{conj:bipartite-comp}: in every $r$-edge-colouring of a bipartite graph with parts of size $m,n$
and appropriate minimum degree, there is a monochromatic component on $\geq \frac{m+n}{r}$ vertices.

\medskip\noindent
\textit{Status:} Proved for $r=2$; weaker bounds for $r \geq 3$.
\end{problem}

\begin{problem}[Partição de $K_{n,n}$ em ciclos]
Is $\mathrm{cp}_2(K_{n,n}) = 3$? (Currently proved $\leq 4$.)
\end{problem}

\begin{problem}[Conjectura de Stein — hipergrafos]
Prove that for any $r, n \in \mathbb{N}$, every $2$-colouring of $K_n^{(r)}$ has a partition
into two monochromatic tight paths of different colours.

\medskip\noindent
\textit{Status:} Proved for $r=2$ (trivial) and $r=3$ (Stein, 2022). Open for $r \geq 4$.
\end{problem}

\begin{problem}[Conjectura de Ryser]\label{prob:ryser}
Prove \Cref{conj:ryser}: $\mathrm{tc}_r(G) \leq (r-1)\alpha(G)$ for every graph $G$.

\medskip\noindent
\textit{Status:} Proved for $r \leq 5$ (when $\alpha(G)=1$, i.e., $G=K_n$).
The case $r=3$ for general graphs was proved by Aharoni. Open for $r \geq 4$ and $\alpha \geq 2$.
\end{problem}

\begin{problem}[Factor $\log r$ em \Cref{thm:dibraccio}]
Is the $\log r$ factor in the upper bound of Di Braccio--Patel (2026) necessary?
Can $\mathrm{cp}_r(\delta) = \Theta\!\left(r \cdot \lceil r/\log(1/\delta)\rceil\right)$ be proved without the $\log r$?
\end{problem}

\begin{problem}[Componentes com muitas arestas para $r \geq 5$]
Determine $\inf_c$ such that every $r$-edge-colouring of $K_n$ contains a monochromatic component
with $\geq c \cdot \binom{n}{2}$ edges, for $r \geq 5$.

\medskip\noindent
\textit{Status:} $r=3 \Rightarrow c=1/6$ (Luo); $r=4 \Rightarrow c=1/12$ (Conlon--Luo--Tyomkyn). Open for $r \geq 5$.
\end{problem}

%%===================================================================
\chapter*{Referências / Bibliography}
\addcontentsline{toc}{chapter}{Referências / Bibliography}
%%===================================================================


\noindent Os 44 artigos desta compilação estão organizados abaixo por identificador arXiv.
Todos os PDFs foram baixados e estão disponíveis na raiz do repositório.


\begin{enumerate}[label={[\arabic*]}, leftmargin=2cm]
\item Pokrovskiy, A.\ (2012). Partitioning edge-coloured complete graphs into monochromatic cycles and paths. \arxiv{1205.5492}

\item Gyárfás, A.\ (2015). Vertex covers by monochromatic pieces — A survey. \arxiv{1509.05539}

\item Gyárfás, A.\ (2017). Vertex covering with monochromatic pieces of few colours. \arxiv{1711.01557}

\item Gyárfás, A.; Király, Z.\ (2016). Covering complete partite hypergraphs by monochromatic components. \arxiv{1604.02791}

\item Grinshpun, A.; Sárközy, G.N.\ (2014). Monochromatic bounded degree subgraph partitions. \arxiv{1405.7507}

\item Gyárfás, A.; Sárközy, G.N.\ (2014). Monochromatic cycle partitions in local edge colourings. \arxiv{1403.5975}

\item Balogh, J.; Barát, J.; Gerbner, D.; Gyárfás, A.; Sárközy, G.N.\ (2015). Partitioning 2-edge-colored graphs by monochromatic paths and cycles. \arxiv{1509.05544}

\item Bal, D.; DeBiasio, L.\ (2015). Partitioning random graphs into monochromatic components. \arxiv{1509.09168}

\item Letzter, S.\ (2015). Monochromatic cycle partitions of 2-coloured graphs with minimum degree $3n/4$. \arxiv{1502.07736}

\item Schaudt, O.; Stein, M.\ (2014). Partitioning two-coloured complete multipartite graphs into monochromatic paths and cycles. \arxiv{1407.5083}

\item Gyárfás, A.; Sárközy, G.N.\ (2015). Local colourings and monochromatic partitions in complete bipartite graphs. \arxiv{1501.05619}

\item Kohayakawa, Y.; Mota, G.O.; Schacht, M.\ (2016). Monochromatic trees in random graphs. \arxiv{1611.10299}

\item Korándi, D.; Sudakov, B.\ (2017). Monochromatic cycle covers in random graphs. \arxiv{1712.03145}

\item Bennett, P.; DeBiasio, L.; Dudek, A.; English, S.\ (2017). Large monochromatic components and long monochromatic cycles in random hypergraphs. \arxiv{1709.02990}

\item DeBiasio, L.; Krueger, R.A.; Sárközy, G.N.\ (2018). Large monochromatic components in multicolored bipartite graphs. \arxiv{1806.05271}

\item Gyárfás, A.; Sárközy, G.N.\ (2018). Long monochromatic paths and cycles in 2-colored bipartite graphs. \arxiv{1806.05119}

\item Korándi, D.; Sárközy, G.N.; Selkow, S.\ (2018). Monochromatic cycle partitions in random graphs. \arxiv{1807.06607}

\item Allen, P.; Böttcher, J.; Lang, R.; Skokan, J.; Stein, M.\ (2019). Minimum degree conditions for monochromatic cycle partitioning. \arxiv{1902.05882}

\item Balogh, J.; Kostochka, A.; Lavrov, M.; Liu, X.\ (2019). Monochromatic connected matchings in 2-edge-colored multipartite graphs. \arxiv{1905.04653}

\item Balogh, J.; Kostochka, A.; Lavrov, M.; Liu, X.\ (2019). Long monochromatic paths and cycles in 2-edge-colored multipartite graphs. \arxiv{1905.04657}

\item Balogh, J.; Kostochka, A.; Lavrov, M.; Liu, X.\ (2019). Monochromatic paths and cycles in 2-edge-colored graphs with large minimum degree. \arxiv{1906.02854}

\item Guggiari, H.; Scott, A.\ (2019). Monochromatic components in edge-coloured graphs with large minimum degree. \arxiv{1909.09178}

\item (2020). Covering 3-edge-coloured random graphs with monochromatic trees. \arxiv{2006.14469}

\item Füredi, Z.; Luo, R.\ (2020). Large monochromatic components in almost complete graphs and bipartite graphs. \arxiv{2008.12217}

\item Bradač, D.; Bucić, M.\ (2021). Covering random graphs with monochromatic trees. \arxiv{2109.02569}

\item Luo, S.\ (2021). On connected components with many edges. \arxiv{2111.13342}

\item Arras, P.\ (2022). Ore- and Pósa-type conditions for partitioning 2-edge-coloured graphs into monochromatic cycles. \arxiv{2202.06388}

\item (2022). Partitioning a 2-edge-coloured graph of minimum degree $2n/3+o(n)$ into three monochromatic cycles. \arxiv{2204.00496}

\item Conlon, D.; Luo, S.; Tyomkyn, M.\ (2022). Monochromatic components with many edges. \arxiv{2204.11360}

\item Stein, M.\ (2022). Monochromatic paths in 2-edge coloured graphs and hypergraphs. \arxiv{2204.12464}

\item Bal, D.; DeBiasio, L.\ (2023). Large monochromatic components in hypergraphs with large minimum codegree. \arxiv{2301.05806}

\item Lichev, L.; Luo, S.\ (2023). Large monochromatic components in colorings of complete hypergraphs. \arxiv{2302.04487}

\item Bal, D.; DeBiasio, L.\ (2023). Large monochromatic components in expansive hypergraphs. \arxiv{2302.06669}

\item Zhang, Y.; Peng, Y.; Luo, Z.\ (2023). Monochromatic cycles in 2-edge-colored bipartite graphs with large minimum degree. \arxiv{2304.08003}

\item (2024). Monochromatic partitions in 2-edge-coloured bipartite graphs. \arxiv{2403.12587}

\item (2024). Partitioning 2-edge-coloured bipartite graphs into monochromatic cycles. \arxiv{2409.03394}

\item Pokrovskiy, A.; Versteegen, L.; Williams, E.\ (2024). A proof of a conjecture of Erdős and Gyárfás on monochromatic path covers. \arxiv{2409.03623}

\item Fox, H.; Luo, S.\ (2025). Monochromatic components with many edges in random graphs. \arxiv{2509.01766}

\item Di Braccio, F.; Patel, V.\ (2026). Monochromatic cycle partitions of $r$-edge-coloured graphs with high minimum degree. \arxiv{2601.22117}

\item Hawranick, L.; Luo, R.\ (2026). Covering complete $r$-partite hypergraphs with few monochromatic components. \arxiv{2603.04704}

\item Girão, A.; Letzter, S.; Sahasrabudhe, J.\ Partitioning a graph into monochromatic connected subgraphs. \arxiv{1708.01284}

\end{enumerate}

\end{document}
